% =====================================================================
\chapter{Модели: создание и основные настройки}
% =====================================================================
\chapterintro
  {как создать, скопировать и выбрать модель, что настраивается на странице Setup: таймеры,
   триммеры, газ, предполётные проверки}
  {пульт после первоначальной настройки}

\section{Что такое «модель» в EdgeTX}

\textbf{Модель} --- это набор всех настроек для одного летательного аппарата (или машины):
входы, миксеры, выходы, таймеры, радиомодуль, телеметрия. В пульте можно хранить десятки
моделей и переключаться между ними. Правило: \textbf{одна физическая модель --- одна модель
в пульте}. Даже если у вас три одинаковых квадрокоптера, лучше сделать три модели (копии) ---
это защищает от случайного взлёта «не той» модели (см. Model Match в главе~\ref{ch:rf}).

\section{Список моделей}

\btn{MDL} открывает меню модели; первая страница --- \menu{Model Select}. Здесь:
\begin{itemize}
  \item колесом выберите модель, \btn{ENTER} или долгое \btn{ENTER} → \menu{Select model} ---
        сделать её текущей;
  \item долгое \btn{ENTER} на пустой строке → \menu{Create model} --- создать новую;
  \item долгое \btn{ENTER} на модели → \menu{Copy model}, \menu{Move model},
        \menu{Delete model}.
\end{itemize}

При создании EdgeTX может предложить мастер (wizard) для типовой модели --- самолёт, крыло,
планер, мультикоптер. Мастер удобен для первого знакомства, но в книге мы будем настраивать
всё вручную, чтобы понимать, что происходит.

\section{Страница Setup}

\btn{PAGE>} от списка моделей ведёт на \mpath{MDL\sep Setup} --- главную страницу модели.

\subsection{Имя модели}

\menu{Model name} --- до 15 символов. Вводится колесом посимвольно; долгое \btn{ENTER} на
символе переключает регистр. Давайте осмысленные имена: \code{Quad-5in}, \code{Wing-900},
\code{Cessna}.

\subsection{Таймеры}

У модели три таймера (\menu{Timer 1}--\menu{Timer 3}). У каждого:

\begin{center}\small
\renewcommand{\arraystretch}{1.2}
\begin{tabularx}{\linewidth}{@{}p{30mm}X@{}}
\toprule
\menu{Mode} & когда идёт отсчёт:
  \menu{OFF} --- выключен; \menu{ON} --- всегда (по переключателю);
  \menu{Start} --- запускается переключателем и больше не останавливается;
  \menu{Throttle} (THs) --- идёт, пока газ не на минимуме;
  \menu{Throttle \%} (TH\%) --- скорость пропорциональна газу;
  \menu{Throttle Start} (THt) --- запускается при первом движении газа и больше не
  останавливается. \\
\menu{Switch} & переключатель, разрешающий счёт (например, арм \sw{SA↓}) \\
\menu{Start} & начальное значение. 0:00 --- счёт вверх, например 4:00 --- обратный отсчёт \\
\menu{Minute call} & голосом объявлять каждую минуту \\
\menu{Countdown} & писк/голос/вибро в последние секунды \\
\menu{Persistent} & сохранять значение между выключениями пульта --- для учёта общего
  налёта модели \\
\bottomrule
\end{tabularx}
\end{center}

\begin{recipe}{таймер полёта для квадрокоптера}
Timer 1: \menu{Mode} = \menu{ON}, \menu{Switch} = переключатель арма (например \sw{SA↓}),
\menu{Start} = \code{4:00}, \menu{Minute call} = вкл., \menu{Countdown} = \menu{Voice}.
Таймер считает время в арме и за 4 минуты голосом напоминает, что пора садиться.
Timer 2: \menu{Mode} = \menu{ON}, тот же переключатель, \menu{Persistent} = \menu{Flight} ---
суммарный налёт модели.
\end{recipe}

\subsection{Триммеры}

\begin{description}[font=\sffamily\bfseries\color{etx}]
  \item[Extended trims] расширенный диапазон триммеров ($\pm100\,\%$ вместо $\pm25\,\%$).
  \item[Trim step] шаг: \menu{Exponential} (мелкий у нуля, крупнее дальше), \menu{Extra fine},
        \menu{Fine}, \menu{Medium}, \menu{Coarse}.
  \item[Display trims] показывать численное значение триммеров на главном экране.
\end{description}

Для квадрокоптеров с полётным контроллером триммеры обычно \emph{не используются}: стики
должны давать ровно 0 в центре, а выравнивание делает контроллер (акселерометр). Триммеры
можно отключить в \menu{Inputs} (параметр \menu{Trim} = \menu{OFF}) или вовсе переназначить
на другие функции.

\subsection{Газ}

\begin{description}[font=\sffamily\bfseries\color{etx}]
  \item[Throttle source] источник, который считается «газом» для таймеров и проверки газа
        при включении. Обычно \sw{Thr}; на планере с электромотором на тумблере ---
        этот тумблер.
  \item[Throttle trim idle only] триммер газа работает только в зоне малого газа ---
        регулирует холостой ход ДВС, не меняя максимум.
  \item[Throttle warning] (в \menu{Preflight checks}) --- проверять газ при включении.
\end{description}

\subsection{Предполётные проверки (Preflight checks)}

\begin{itemize}
  \item \menu{Throttle state} --- проверять, что газ на минимуме.
  \item \menu{Custom position} --- вместо минимума задать свою позицию газа.
  \item \menu{Switch positions} --- список переключателей с требуемыми положениями.
        Встаньте на этот пункт, выставьте переключатели «как для старта» и долгим \btn{ENTER}
        запомните. Лишние переключатели можно исключить из проверки.
  \item \menu{Pots / Sliders} --- аналогично для крутилок.
\end{itemize}

\begin{tip}
Включайте проверку переключателей хотя бы для переключателя арма. Тогда пульт не
«забудет» заармленный квадрокоптер при включении и не даст случайно запустить мотор.
\end{tip}

\subsection{Прочие пункты}

\begin{itemize}
  \item \menu{Center beep} --- писк при проходе стиком или крутилкой через центр.
  \item \menu{Use global functions} --- применять глобальные функции пульта к этой модели.
  \item \menu{Internal RF}, \menu{External RF} --- радиомодули, глава~\ref{ch:rf}.
  \item \menu{Trainer} --- режим тренера для этой модели, глава~\ref{ch:trainer}.
  \item \menu{USB Joystick} --- как модель выглядит для компьютера в режиме джойстика.
\end{itemize}

\begin{summary}
\begin{itemize}
  \item Каждая физическая модель --- отдельная модель в пульте.
  \item \menu{Setup}: имя, таймеры, триммеры, источник газа, предполётные проверки, радиомодуль.
  \item Предполётные проверки --- дешёвая и эффективная защита от неприятностей.
\end{itemize}
\end{summary}
